%==============================================================================
%  Lab Report Template — Lab 2: Supervised Learning
%  Course : Introduction to Machine Learning (with Python & scikit-learn)
%  Follows: notebooks/02-supervised-learning-code.ipynb
%           "Comparing Classifiers on the Breast Cancer Dataset"
%
%  Compile: pdflatex lab2-supervised-learning.tex   (run twice for references)
%       or: latexmk -pdf lab2-supervised-learning.tex
%
%  HOW TO USE
%  ----------
%  1. Fill in the metadata in the "STUDENT INFORMATION" block below.
%  2. Run the notebook, then replace every  % TODO  / \textit{TODO} with the
%     numbers and plots you obtained.
%  3. Put your figures in a  figures/  subfolder and use \includegraphics.
%  4. Delete the guidance comments (lines starting with %) before submitting.
%==============================================================================
\documentclass[11pt,a4paper]{article}

%------------------------------ packages --------------------------------------
\usepackage[T1]{fontenc}
\usepackage[utf8]{inputenc}
\usepackage{lmodern}
\usepackage[margin=2.5cm]{geometry}
\usepackage{graphicx}
\usepackage{booktabs}
\usepackage{amsmath,amssymb}
\usepackage{xcolor}
\usepackage{listings}
\usepackage{caption}
\usepackage{subcaption}
\usepackage{enumitem}
\usepackage{float}
\usepackage{fancyhdr}
\usepackage[hidelinks]{hyperref}

%------------------------------ metadata --------------------------------------
% ========================= STUDENT INFORMATION ===============================
\newcommand{\Course}{Introduction to Machine Learning}
\newcommand{\LabNo}{Lab 2}
\newcommand{\LabTitle}{Supervised Learning}
\newcommand{\StudentName}{Your Name}          % TODO: your full name
\newcommand{\StudentID}{Student ID}           % TODO: your student number
\newcommand{\DueDate}{2026.10.21}                 % TODO: the due date
\newcommand{\Instructor}{Zhang Li \& Jing Ming}
\newcommand{\Institution}{Shandong Women's University}
% =============================================================================

%------------------------------ colours & code style --------------------------
\definecolor{codebg}{RGB}{245,242,247}
\definecolor{kw}{RGB}{125,15,98}      % keywords (burgundy)
\definecolor{str}{RGB}{31,51,82}      % strings (navy)
\definecolor{cmt}{RGB}{91,100,114}    % comments (muted)
\definecolor{accent}{RGB}{138,28,48}  % headings accent

\lstdefinestyle{py}{
  language=Python,
  backgroundcolor=\color{codebg},
  basicstyle=\ttfamily\small,
  keywordstyle=\color{kw}\bfseries,
  stringstyle=\color{str},
  commentstyle=\color{cmt}\itshape,
  numberstyle=\tiny\color{cmt},
  numbers=left,
  numbersep=8pt,
  frame=leftline,
  framerule=1.2pt,
  rulecolor=\color{accent},
  showstringspaces=false,
  breaklines=true,
  tabsize=4,
  xleftmargin=14pt,
  aboveskip=10pt,
  belowskip=6pt,
}
\lstset{style=py}

%------------------------------ headers/footers -------------------------------
\pagestyle{fancy}
\fancyhf{}
\lhead{\small \Course}
\rhead{\small \LabNo\ — \LabTitle}
\cfoot{\small \StudentName\ (\StudentID)}
\rfoot{\small \thepage}
\renewcommand{\headrulewidth}{0.4pt}
\renewcommand{\footrulewidth}{0.4pt}

%------------------------------ helpers ---------------------------------------
% A grey placeholder box to stand in for a figure you have not added yet.
\newcommand{\placeholder}[2][8cm]{%
  \fbox{\parbox[c][5cm][c]{#1}{\centering\color{cmt}\itshape #2}}%
}

%==============================================================================
\begin{document}

%------------------------------ title block -----------------------------------
\begin{titlepage}
\centering
\vspace*{2cm}
{\LARGE\bfseries \Course\par}
\vspace{0.6cm}
{\color{accent}\rule{\linewidth}{1.5pt}}\par
\vspace{0.5cm}
{\Huge\bfseries \LabNo: \LabTitle\par}
\vspace{0.3cm}
{\large Comparing Classifiers on the Breast Cancer Dataset\par}
\vspace{0.4cm}
{\color{accent}\rule{\linewidth}{1.5pt}}\par
\vspace{2.3cm}

\begin{tabular}{rl}
  \textbf{Student:}    & \StudentName \\[4pt]
  \textbf{Student ID:} & \StudentID \\[4pt]
  \textbf{Instructor:} & \Instructor \\[4pt]
  \textbf{Institution:}& \Institution \\[4pt]
  \textbf{Date:}       & \DueDate \\
\end{tabular}

\vfill
{\color{cmt}\small Based on \texttt{notebooks/02-supervised-learning-code.ipynb}.
Replace all placeholder text before submission.}
\end{titlepage}

\tableofcontents
\newpage

%==============================================================================
\section{Objectives}
%------------------------------------------------------------------------------
In this lab you compare several supervised classifiers on the same dataset.
By the end you should be able to:
\begin{itemize}[leftmargin=*]
  \item Train several classifiers and compare \textbf{train vs.\ test} accuracy.
  \item Observe \textbf{overfitting} in an unrestricted decision tree.
  \item Read \textbf{feature importances} from a random forest.
  \item Explain why \textbf{SVMs need feature scaling}.
  \item Obtain class \textbf{probability estimates} with \lstinline|predict_proba|.
\end{itemize}

%==============================================================================
\section{Background}
%------------------------------------------------------------------------------
In \emph{supervised learning} a model is trained on labelled examples
$(\mathbf{x}_i, y_i)$ and learns a mapping $f:\mathbf{x}\mapsto y$ that should
\textbf{generalise} to unseen data. Here $y\in\{0,1\}$ (malignant / benign), so
this is a \textbf{binary classification} task. We report accuracy on a held-out
test set,
\[
  \text{accuracy} = \frac{\text{\# correct predictions}}{\text{\# samples}},
\]
and compare it to the training accuracy to detect over-/under-fitting:
\begin{itemize}[leftmargin=*]
  \item \textbf{Overfitting} — high training accuracy but a noticeably lower test
        accuracy (the model memorised training quirks).
  \item \textbf{Underfitting} — low accuracy on both sets (the model is too simple).
\end{itemize}

\noindent This lab touches five model families: \textbf{logistic regression},
\textbf{decision trees}, \textbf{random forests}, \textbf{support vector machines
(SVM)}, and uses \textbf{feature scaling} (\lstinline|StandardScaler|) and
\textbf{probability estimates} (\lstinline|predict_proba|).

%==============================================================================
\section{Dataset}
%------------------------------------------------------------------------------
We use the \textbf{Breast Cancer Wisconsin} dataset, bundled with scikit-learn
as \lstinline|load_breast_cancer()|. Each sample describes a tumour with 30
numeric features; the target labels it \emph{malignant} or \emph{benign}.

\begin{table}[H]
\centering
\caption{Dataset summary (fixed facts, provided for reference).}
\label{tab:dataset}
\begin{tabular}{@{}lr@{}}
\toprule
\textbf{Property} & \textbf{Value} \\
\midrule
Source              & \lstinline|sklearn.datasets.load_breast_cancer| \\
Task type           & Binary classification \\
Number of samples   & 569 \\
Number of features  & 30 (numeric) \\
Classes             & malignant (0), benign (1) \\
Split               & \lstinline|train_test_split(..., random_state=0)| \\
Training samples    & 426 \\
Test samples        & 143 \\
\bottomrule
\end{tabular}
\end{table}

%==============================================================================
\section{Methodology}
%------------------------------------------------------------------------------
The notebook is organised as six exercises. All models come from scikit-learn
and share the same train/test split (\lstinline|random_state=0|). The setup:

\begin{lstlisting}[caption={Imports and data split (Exercise 1).},label={lst:setup}]
import numpy as np
from sklearn.datasets import load_breast_cancer
from sklearn.model_selection import train_test_split
from sklearn.linear_model import LogisticRegression
from sklearn.tree import DecisionTreeClassifier
from sklearn.ensemble import RandomForestClassifier
from sklearn.svm import SVC
from sklearn.preprocessing import StandardScaler

cancer = load_breast_cancer()
X_train, X_test, y_train, y_test = train_test_split(
    cancer.data, cancer.target, random_state=0)   # -> 426 train / 143 test
\end{lstlisting}

\noindent The six exercises are:
\begin{enumerate}[leftmargin=*]
  \item Load the data and create the train/test split.
  \item Train \lstinline|LogisticRegression(max_iter=5000)|; report train/test accuracy.
  \item Train a \lstinline|DecisionTreeClassifier| with no depth limit and with
        \lstinline|max_depth=4|; compare train vs.\ test to see overfitting.
  \item Train \lstinline|RandomForestClassifier(n_estimators=100, random_state=0)|;
        report test accuracy and the 5 most important features.
  \item Train \lstinline|SVC()| on raw vs.\ \lstinline|StandardScaler|-scaled data.
  \item Print \lstinline|predict_proba| for the first 5 test samples.
\end{enumerate}

%==============================================================================
\section{Experiments and Results}
%------------------------------------------------------------------------------
% Run the notebook and fill in every \textit{TODO} with the value you obtained.

\subsection{Exercise 1 — Data split}
Report the sizes produced by the split.
\begin{itemize}[leftmargin=*]
  \item Training samples: \textit{TODO} \quad(expected 426)
  \item Test samples: \textit{TODO} \quad(expected 143)
\end{itemize}

\subsection{Exercise 2 — Logistic regression}
\begin{table}[H]
\centering
\caption{Logistic regression accuracy.}
\label{tab:logreg}
\begin{tabular}{@{}lcc@{}}
\toprule
\textbf{Model} & \textbf{Train acc.} & \textbf{Test acc.} \\
\midrule
\lstinline|LogisticRegression(max_iter=5000)| & \textit{TODO} & \textit{TODO} \\
\bottomrule
\end{tabular}
\end{table}

\subsection{Exercise 3 — Decision tree and overfitting}
Compare an unrestricted tree with a depth-limited one.
\begin{table}[H]
\centering
\caption{Decision tree: effect of \texttt{max\_depth} on over-fitting.}
\label{tab:tree}
\begin{tabular}{@{}lccl@{}}
\toprule
\textbf{Model} & \textbf{Train acc.} & \textbf{Test acc.} & \textbf{Comment} \\
\midrule
\lstinline|DecisionTree(random_state=0)|          & \textit{TODO} & \textit{TODO} & overfits? \\
\lstinline|DecisionTree(max_depth=4)|             & \textit{TODO} & \textit{TODO} & \\
\bottomrule
\end{tabular}
\end{table}
% Optional extension: plot train/test accuracy against max_depth = 1..10.
\begin{figure}[H]
\centering
% \includegraphics[width=0.7\linewidth]{figures/tree_depth_curve.png}
\placeholder[10cm]{(Optional) Plot of train/test accuracy\\ vs.\ \texttt{max\_depth} = 1..10.}
\caption{Train/test accuracy as the tree depth increases (optional extension).}
\label{fig:depth}
\end{figure}

\subsection{Exercise 4 — Random forest and feature importances}
Report the random-forest test accuracy and list the five most important features.
\begin{table}[H]
\centering
\caption{Top-5 features by random-forest importance.}
\label{tab:importance}
\begin{tabular}{@{}clr@{}}
\toprule
\textbf{Rank} & \textbf{Feature name} & \textbf{Importance} \\
\midrule
1 & \textit{TODO} & \textit{TODO} \\
2 & \textit{TODO} & \textit{TODO} \\
3 & \textit{TODO} & \textit{TODO} \\
4 & \textit{TODO} & \textit{TODO} \\
5 & \textit{TODO} & \textit{TODO} \\
\bottomrule
\end{tabular}
\end{table}
\noindent Random-forest test accuracy: \textit{TODO} \quad(expected $>0.90$).

\begin{figure}[H]
\centering
% \includegraphics[width=0.75\linewidth]{figures/feature_importances.png}
\placeholder[11cm]{(Optional) Bar chart of \texttt{rf.feature\_importances\_}.}
\caption{Random-forest feature importances (optional).}
\label{fig:importance}
\end{figure}

\subsection{Exercise 5 — Scaling matters for SVMs}
Compare the SVM on raw features versus \lstinline|StandardScaler|-scaled features.
\begin{table}[H]
\centering
\caption{SVM test accuracy with and without feature scaling.}
\label{tab:svm}
\begin{tabular}{@{}lc@{}}
\toprule
\textbf{Preprocessing} & \textbf{Test acc.} \\
\midrule
\lstinline|SVC()| on raw data          & \textit{TODO} \\
\lstinline|SVC()| on scaled data       & \textit{TODO} \\
\bottomrule
\end{tabular}
\end{table}

\subsection{Exercise 6 — Probability estimates}
Report \lstinline|predict_proba| for the first five test samples.
\begin{table}[H]
\centering
\caption{Predicted class probabilities for the first 5 test samples.}
\label{tab:proba}
\begin{tabular}{@{}cccl@{}}
\toprule
\textbf{Sample} & \textbf{P(malignant)} & \textbf{P(benign)} & \textbf{Predicted} \\
\midrule
1 & \textit{TODO} & \textit{TODO} & \textit{TODO} \\
2 & \textit{TODO} & \textit{TODO} & \textit{TODO} \\
3 & \textit{TODO} & \textit{TODO} & \textit{TODO} \\
4 & \textit{TODO} & \textit{TODO} & \textit{TODO} \\
5 & \textit{TODO} & \textit{TODO} & \textit{TODO} \\
\bottomrule
\end{tabular}
\end{table}

%==============================================================================
\section{Discussion}
%------------------------------------------------------------------------------
Answer with evidence from your tables and figures above:
\begin{itemize}[leftmargin=*]
  \item \textbf{Overfitting (Ex.~3).} How large is the gap between train and test
        accuracy for the unrestricted tree? Does \lstinline|max_depth=4| reduce it?
  \item \textbf{Best model.} Which classifier generalised best on the test set, and
        why do you think so?
  \item \textbf{Scaling (Ex.~5).} How much did \lstinline|StandardScaler| change the
        SVM's test accuracy? Why are SVMs sensitive to feature scales while trees
        and forests are not?
  \item \textbf{Feature importances (Ex.~4).} Do the top features in
        Table~\ref{tab:importance} make intuitive sense for tumour diagnosis?
  \item \textbf{Probabilities (Ex.~6).} Are any of the five samples predicted with
        low confidence? How could a probability threshold be useful in a medical
        setting?
\end{itemize}
% TODO: write 1--2 paragraphs answering the questions above.

%==============================================================================
\section{Conclusion}
%------------------------------------------------------------------------------
% 3--5 sentences summarising what you learned and the main result.
% TODO.

%==============================================================================
% \begin{thebibliography}{9}
% \bibitem{mueller2016}
% A.~C. M\"uller and S.~Guido,
% \emph{Introduction to Machine Learning with Python}.
% O'Reilly Media, 2016.

% \bibitem{sklearn}
% F.~Pedregosa et al.,
% ``Scikit-learn: Machine Learning in Python,''
% \emph{Journal of Machine Learning Research}, vol.~12, pp.~2825--2830, 2011.
% \end{thebibliography}

%==============================================================================
\appendix
\section{Source Code}
%------------------------------------------------------------------------------
% Paste your completed code, or reference the notebook you ran.
Full code: \texttt{notebooks/02-supervised-learning-code.ipynb}. Key solution:

\begin{lstlisting}[caption={Reference solution for Exercises 2--6.}]
# Ex2 - logistic regression
logreg = LogisticRegression(max_iter=5000).fit(X_train, y_train)
print("logreg train/test:", logreg.score(X_train, y_train),
      logreg.score(X_test, y_test))

# Ex3 - decision tree (overfitting)
full   = DecisionTreeClassifier(random_state=0).fit(X_train, y_train)
pruned = DecisionTreeClassifier(max_depth=4, random_state=0).fit(X_train, y_train)

#Optional: Train/test accuracy as the tree depth increases (optional extension).  Plot of train/test accuracy vs. max_depth = 1..10

import matplotlib.pyplot as plt

plt.figure(figsize=(8, 5))
train_acc = []
test_acc = []
for depth in range(1, 11):
    tree = DecisionTreeClassifier(max_depth=depth, random_state=0).fit(X_train, y_train)
    train_acc.append(tree.score(X_train, y_train))
    test_acc.append(tree.score(X_test, y_test))
plt.plot(range(1, 11), train_acc, label='Train Accuracy', marker='o')
plt.plot(range(1, 11), test_acc, label='Test Accuracy', marker='o')
plt.xlabel('Max Depth')
plt.ylabel('Accuracy')
plt.title('Decision Tree Accuracy vs Max Depth')
plt.xticks(range(1, 11))
plt.legend()
plt.grid()
plt.show()

# Ex4 - random forest + feature importances
rf = RandomForestClassifier(n_estimators=100, random_state=0).fit(X_train, y_train)
top = np.argsort(rf.feature_importances_)[::-1][:5]
print(rf.score(X_test, y_test), list(cancer.feature_names[top]))

# Optional. Random-forest feature importances (optional).  Bar chart of rf.feature_importances_

plt.figure(figsize=(10, 6))
importances = rf.feature_importances_
indices = np.argsort(importances)[::-1]
plt.title('Feature Importances from Random Forest')
plt.bar(range(X_train.shape[1]), importances[indices], align='center')
plt.xticks(range(X_train.shape[1]), cancer.feature_names[indices], rotation=90)
plt.tight_layout()
plt.show()

# Ex5 - scaling matters for SVM
scaler = StandardScaler().fit(X_train)
svc_raw    = SVC().fit(X_train, y_train)
svc_scaled = SVC().fit(scaler.transform(X_train), y_train)

# Ex6 - probability estimates
print(np.round(rf.predict_proba(X_test[:5]), 3))
\end{lstlisting}

\end{document}
